\subsection{Ordered Fields, Suprema, and Real Structure}

% **Role:** rigorous language of real and complex analysis.
%
% Natural numbers, integers, rationals, reals, complex numbers, ordered fields, supremum property, Archimedean property, density.


Number systems supply the ordered and algebraic background for Analysis: $\Q$ has gaps, $\R$ repairs them by completeness, and $\C$ keeps field algebra without order.

\begin{definition} \label{def:inductive-natural-rational-and-dense-subsets} \hlt{Natural, Rational, Irrational, and Dense Subsets}\\
A set $E\subseteq\R$ is \hlt{inductive} iff $1\in E$ and, for all $x\in\R$, $x\in E$ implies $x+1\in E$.\\
Let $\mathcal{I}$ be the family of inductive subsets of $\R$.
\begin{enumerate}[label=(\roman*)]
\setlength{\itemsep}{0pt}
\item $\N\defeq\bigcap_{E\in\mathcal{I}}E$, the least inductive subset of $\R$.
\item $\Z\defeq\N\cup\{0\}\cup\{-n:n\in\N\}$.
\item $\Q\defeq\{m/n:m,n\in\Z,\ n\neq0\}$; $x\in\R$ is \hlt{irrational} iff $x\notin\Q$.
\item $D\subseteq\R$ is \hlt{dense in $\R$} iff for all $a,b\in\R$ with $a<b$, there exists $d\in D$ with $a<d<b$.
\end{enumerate}
\end{definition}

\begin{definition} \label{def:field-and-ordered-field-analysis} \hlt{Fields and Ordered Fields}\\
Let $F$ be a set with operations $+,\cdot:F\times F\to F$ and elements $0,1\in F$ with $0\neq1$.
\begin{enumerate}[label=(\roman*)]
\setlength{\itemsep}{0pt}
\item A \hlt{field} has $(F,+,0)$ and $(F\setminus\{0\},\cdot,1)$ abelian groups and $x(y+z)=xy+xz$ for all $x,y,z\in F$.
\item An \hlt{ordered field} is a field with total order $<$ satisfying $y<z\Rightarrow x+y<x+z$ and $0<x,0<y\Rightarrow0<xy$.
\end{enumerate}
\end{definition}

\begin{principle} \label{princ:completeness-axiom-real-line} \hlt{Completeness Axiom}\\
Every nonempty set $E\subseteq\R$ that is bounded above has a least upper bound in $\R$, denoted $\sup E$.\\
Equivalently, the real line has no missing cut between a nonempty bounded set and its upper bounds.
\end{principle}

\begin{remark} \label{rem:fields-as-analysis-objects-and-linear-scalars} \hlt{Fields as Objects and Scalars}\\
In analysis, $\Q$ and $\R$ are ordered fields: $\Q$ lacks the least-upper-bound property, while $\R$ has it.
\end{remark}

\begin{proposition} \label{prop:field-axiom-consequences-analysis} \hlt{Basic Consequences of the Field Axioms}\\
Let $F$ be a field. For all $a,b,c,x,y\in F$:
\begin{enumerate}[label=(\roman*)]
\setlength{\itemsep}{0pt}
\item $a+c=b+c$ implies $a=b$; if $c\neq0$, then $ac=bc$ implies $a=b$.
\item The identities $0,1$ and inverses $-a$, $a^{-1}$ $(a\neq0)$ are unique.
\item $0x=0$.
\item $xy=0$ implies $x=0$ or $y=0$.
\item $(-x)y=-(xy)=x(-y)$ and $(-x)(-y)=xy$.
\end{enumerate}
\begin{proof}
\begin{enumerate}[label=(\roman*)]
\setlength{\itemsep}{0pt}
\item Add $-c$ to $a+c=b+c$. If $c\neq0$, multiply $ac=bc$ by $c^{-1}$.
\item In any group $(G,*,e)$, if $e'$ is an identity, $e=e*e'=e'$.\\
If $u*a=e=a*v$, then $u=u*e=u*(a*v)=(u*a)*v=v$. Apply to $(F,+,0)$ and $(F\setminus\{0\},\cdot,1)$.
\item Since $0x+0x=(0+0)x=0x$, cancellation gives $0x=0$.
\item If $xy=0$ and $x\neq0$, multiply by $x^{-1}$ to get $y=0$.
\item $xy+(-x)y=(x+(-x))y=0$, so $(-x)y=-(xy)$; $x(-y)=-(xy)$. Also $(-x)(-y)=-(x(-y))=xy$.
\end{enumerate}
\hlt{Structural Commentary}: Inverses cancel terms.\\
\hlt{Key Insight}: Negatives are unique additive inverses.
\end{proof}
\end{proposition}

\begin{proposition} \label{prop:ordered-field-inequality-rules-analysis} \hlt{Basic Inequality Rules in Ordered Fields}\\
Let $F$ be an ordered field. For all $c,x,y\in F$:
\begin{enumerate}[label=(\roman*)]
\setlength{\itemsep}{0pt}
\item $x>0$ iff $-x<0$.
\item If $c>0$ and $x<y$, then $cx<cy$.
\item If $c<0$ and $x<y$, then $cy<cx$.
\item If $x\neq0$, then $x^2>0$; in particular, $1>0$.
\item If $0<x<y$, then $0<1/y<1/x$.
\end{enumerate}
\begin{proof}
\begin{enumerate}[label=(\roman*)]
\setlength{\itemsep}{0pt}
\item If $x>0$, add $-x$ to $0<x$ to get $-x<0$. If $-x<0$, add $x$ to get $0<x$.
\item Since $x<y$, $y-x>0$. With $c>0$, $c(y-x)>0$, so $cy-cx>0$ and $cx<cy$.
\item Since $c<0$, (i) gives $-c>0$. By (ii), $(-c)x<(-c)y$; hence $-cx<-cy$, so $cy<cx$.
\item If $x>0$, then $x^2>0$. If $x<0$, then $-x>0$ by (i), so $(-x)^2>0$ and $x^2=(-x)^2$.\\
Since $x\neq0$, totality gives one of these cases. Taking $x=1$ gives $1=1^2>0$.
\item By (iv), $(1/x)^2>0$ and $(1/y)^2>0$. Since $x,y>0$, $1/x=x(1/x)^2>0$ and $1/y=y(1/y)^2>0$, so $(1/x)(1/y)>0$. Apply (ii) to $x<y$ with multiplier $(1/x)(1/y)$ to get $1/y<1/x$.
\end{enumerate}
\hlt{Structural Commentary}: Positivity preserves order under multiplication.\\
\hlt{Key Insight}: Negative multiplication is positive multiplication by $-c$.
\end{proof}
\end{proposition}

\begin{proposition} \label{prop:rational-square-root-two-gap} \hlt{The Rational Gap at $\sqrt{2}$}\\
There is no rational $p$ such that $p^2=2$.
\begin{proof} Let $p=m/n$. Then $m^2=2n^2$, so $m$ is even; $m=2k$. Hence $n^2=2k^2$, so $n$ is even, a contradiction.\\
\hlt{Structural Commentary}: This exhibits a gap in $\Q$.\\
\hlt{Key Insight}: $\sqrt{2}$ is not rational.
\end{proof}
\end{proposition}

\begin{figure}[H]
  \centering
  \begin{tikzpicture}[
  scale=0.95,
  every node/.style={font=\small},
  side/.style={gray!70, line width=0.55pt},
  diagonal/.style={blue!70!black, line width=0.75pt},
  point/.style={circle, fill=gray!85, inner sep=1.15pt},
  labelbox/.style={fill=white, fill opacity=0.88, text opacity=1, inner sep=1.1pt}
]
  \coordinate (A) at (0,0);
  \coordinate (B) at (1.3,0);
  \coordinate (C) at (0,1.3);
  \coordinate (D) at (1.3,1.3);

  \draw[side] (A) -- (B) -- (D) -- (C) -- cycle;
  \draw[diagonal] (A) -- (D);

  \node[point] at (A) {};
  \node[point] at (B) {};
  \node[point] at (C) {};
  \node[point] at (D) {};

  \node[below left] at (A) {$A$};
  \node[below right] at (B) {$B$};
  \node[above left] at (C) {$C$};
  \node[above right] at (D) {$D$};

  \node[below] at ($(A)!0.5!(B)$) {$1$};
  \node[right] at ($(B)!0.5!(D)$) {$1$};
  \node[labelbox, rotate=45] at ($(A)!0.52!(D)$) {$\sqrt{2}$};
\end{tikzpicture}

  \caption{The diagonal of a unit square has length $\sqrt{2}$, so the rational gap is geometric as well as algebraic.}
  \label{fig:unit-square-sqrt-two}
\end{figure}

\begin{example} \label{ex:rational-cut-around-square-root-two} \hlt{No Largest Lower Side and No Smallest Upper Side}\\
Let $A\defeq\{p\in\Q:p>0,\ p^2<2\}$, and $B\defeq\{p\in\Q:p>0,\ p^2>2\}$.\\
For every $p>0$ define $q\defeq p-\frac{p^2-2}{p+2}=\frac{2p+2}{p+2}$. Then $q^2-2=\frac{2(p^2-2)}{(p+2)^2}$.
Hence if $p\in A$, then $q\in A$ and $q>p$; if $p\in B$, then $q\in B$ and $0<q<p$. Thus $A$ has no largest element in $\Q$, and $B$ has no smallest element in $\Q$.
\end{example}

\begin{remark} \label{rem:rational-density-does-not-fill-gaps} \hlt{Density Does Not Mean Completeness}\\
Between two rationals $r<s$ there is another rational, namely $(r+s)/2$. Nevertheless, $\Q$ is not complete.
\end{remark}

\begin{figure}[H]
  \centering
  \begin{tikzpicture}[
  >=Stealth,
  every node/.style={font=\small},
  axis/.style={->, gray!70, line width=0.55pt},
  setbar/.style={blue!70!black, line width=2.0pt, line cap=round},
  negbar/.style={orange!80!black, line width=2.0pt, line cap=round},
  tick/.style={gray!70, line width=0.55pt},
  guide/.style={gray!55, dashed, line width=0.45pt},
  infguide/.style={green!45!black, dashed, line width=0.55pt},
  mguide/.style={purple!70!black, dashed, line width=0.55pt},
  zeroarrow/.style={->, teal!60!black, line width=0.55pt},
  marklabel/.style={fill=white, inner sep=1pt}
]
  \draw[axis] (-4.5,1.25) -- (4.5,1.25);
  \draw[axis] (-4.5,-1.25) -- (4.5,-1.25);

  \draw[teal!60!black, line width=0.55pt] (0,1.14) -- (0,1.36);
  \draw[teal!60!black, line width=0.55pt] (0,-1.36) -- (0,-1.14);
  \node[below=3pt, teal!60!black, marklabel] at (0,1.25) {$0$};
  \node[above=3pt, teal!60!black, marklabel] at (0,-1.25) {$0$};

  \coordinate (infB) at (-1.75,1.25);
  \coordinate (supNegB) at (1.75,-1.25);
  \coordinate (m) at (-3.05,1.25);
  \coordinate (negm) at (3.05,-1.25);

  \draw[infguide] (infB) -- (supNegB);
  \draw[mguide] (m) -- (negm);

  \draw[setbar] (-1.75,1.25) -- (3.7,1.25);
  \foreach \x in {-1.32,-0.75,0.2,1.25,2.2,3.1} {
    \fill[blue!70!black] (\x,1.25) circle (1.2pt);
  }
  \node[above=6pt, blue!70!black, marklabel] at (2.7,1.25) {$B$};

  \draw[negbar] (-3.7,-1.25) -- (1.75,-1.25);
  \foreach \x in {-3.1,-2.2,-1.25,-0.2,0.75,1.32} {
    \fill[orange!80!black] (\x,-1.25) circle (1.2pt);
  }
  \node[below=6pt, orange!80!black, marklabel] at (-2.7,-1.25) {$-B$};

  \draw[tick] (infB) ++(0,-0.11) -- ++(0,0.22);
  \draw[tick] (supNegB) ++(0,-0.11) -- ++(0,0.22);
  \draw[tick] (m) ++(0,-0.11) -- ++(0,0.22);
  \draw[tick] (negm) ++(0,-0.11) -- ++(0,0.22);

  \node[below=6pt, green!45!black, marklabel] at (infB) {$\inf B$};
  \node[below=6pt, green!45!black, marklabel] at (supNegB) {$\sup(-B)$};
  \node[above=6pt, purple!70!black, marklabel] at (m) {$m$};
  \node[below=6pt, purple!70!black, marklabel] at (negm) {$-m$};

  \draw[zeroarrow] (0,0.9) -- (0,-0.9);
  \node[right=4pt, teal!60!black, marklabel] at (0,0) {$x\mapsto -x$};
\end{tikzpicture}

  \caption{Reflecting $B$ through $0$ turns lower bounds of $B$ into upper bounds of $-B$, so $\inf B=-\sup(-B)$.}
  \label{fig:infimum-reflection-via-supremum}
\end{figure}

\begin{corollary} \label{cor:infimum-as-negative-supremum} \hlt{Infimum by Reflection}\\
Let $B\subseteq\R$ be nonempty and bounded below, and define $-B\defeq\{-b:b\in B\}$.\\
Then $-B$ is nonempty and bounded above, and $\inf B=-\sup(-B)$.
\begin{proof} If $m$ is a lower bound of $B$, then $-m$ is an upper bound of $-B$; hence $s\defeq\sup(-B)$ exists.\\
For $b\in B$, $-b\leq s$, so $-s\leq b$; thus $-s$ is a lower bound of $B$.\\
If $t$ is any lower bound of $B$, then $-t$ is an upper bound of $-B$, so $s\leq -t$ and $t\leq -s$. Therefore $-s=\inf B$.\\
\hlt{Structural Commentary}: Negation reverses bounds.\\
\hlt{Key Insight}: $\inf B=-\sup(-B)$.
\end{proof}
\end{corollary}


\begin{definition} \label{def:real-absolute-value-and-distance} \hlt{Real Absolute Value and Distance}\\
For $x\in\R$, define $\abs{x}\defeq
\begin{cases}
x, & x\geq0,\\
-x, & x<0.
\end{cases}$.\\
For $a,b\in\R$, the distance between $a$ and $b$ on the real line is $\abs{a-b}$.
\end{definition}

\begin{proposition} \label{prop:real-absolute-value-triangle-inequality} \hlt{Real Absolute-Value Laws}\\
For $x,y,a,b,c\in\R$,
\begin{enumerate}[label=(\roman*)]
\setlength{\itemsep}{0pt}
\item $\abs{x}\geq0$, $\abs{x}=0$ iff $x=0$, and $-\abs{x}\leq x\leq\abs{x}$;
\item $\abs{xy}=\abs{x}\abs{y}$;
\item $\abs{x+y}\leq\abs{x}+\abs{y}$;
\item $\abs{a-b}\leq\abs{a-c}+\abs{c-b}$.
\end{enumerate}
\begin{proof}
\begin{enumerate}[label=(\roman*)]
\setlength{\itemsep}{0pt}
\item If $x\geq0$, then $\abs{x}=x$; if $x<0$, then $\abs{x}=-x>0$. Hence $\abs{x}\geq0$, and $\abs{x}=0$ iff $x=0$.\\
The bounds $-\abs{x}\leq x\leq\abs{x}$ follow from the same two cases.
\item If $x,y$ have the same sign, then $xy\geq0$ and $\abs{xy}=xy=\abs{x}\abs{y}$.\\
If they have opposite signs, then $xy<0$ and $\abs{xy}=-xy=\abs{x}\abs{y}$.
\item By (i), $-\abs{x}\leq x\leq\abs{x}$ and $-\abs{y}\leq y\leq\abs{y}$, so $-(\abs{x}+\abs{y})\leq x+y\leq\abs{x}+\abs{y}$.\\
Since $\abs{x}+\abs{y}\geq0$, (i) applied to $x+y$ gives $\abs{x+y}\leq\abs{x}+\abs{y}$.
\item Apply (iii) to $a-b=(a-c)+(c-b)$.
\end{enumerate}
\hlt{Structural Commentary}: Absolute value is two-sided order control.\\
\hlt{Key Insight}: Apply (iii) to $(a-c)+(c-b)$.
\end{proof}
\end{proposition}

\begin{perspective} \label{per:dedekind-construction-as-completeness-proof} \hlt{Completeness Becomes a Theorem}\\
Dedekind cuts build completeness from $\Q$: real numbers are rational cuts, rational cuts embed $\Q$, and least upper bounds arise by taking unions of cuts.
\end{perspective}

\begin{definition} \label{def:dedekind-cuts-and-order} \hlt{Dedekind Cuts and Their Order}\\
A \hlt{cut} is a set $\alpha\subseteq\Q$ such that
\begin{enumerate}[label=(\roman*)]
\setlength{\itemsep}{0pt}
\item $\alpha$ is nonempty and $\alpha\neq\Q$;
\item if $p\in\alpha$, $q\in\Q$, and $q<p$, then $q\in\alpha$;
\item if $p\in\alpha$, then some $r\in\alpha$ satisfies $p<r$.
\end{enumerate}
Let $\mathscr{D}$ be the set of all cuts. For $\alpha,\beta\in\mathscr{D}$, define $\alpha<\beta$ to mean $\alpha\subsetneq\beta$.\\
Rational letters $p,q,r,\ldots$ denote elements of $\Q$; Greek letters $\alpha,\beta,\gamma,\ldots$ denote cuts.\\
A cut is the rational left side of a real number: a proper initial segment with no last point.
\end{definition}

\begin{figure}[H]
  \centering
  \begin{tikzpicture}[
  >=Stealth,
  every node/.style={font=\small},
  axis/.style={->, gray!70, line width=0.5pt},
  leftcut/.style={blue!60!black, line width=2.0pt},
  rightpart/.style={orange!75!black, line width=2.0pt},
  marker/.style={gray!55, dashed, line width=0.45pt},
  qpoint/.style={circle, fill=black!65, inner sep=1.05pt}
]
  \draw[axis] (-4.25,0) -- (4.55,0) node[right] {$\Q$};
  \draw[leftcut] (-4.0,0) -- (-0.08,0);
  \draw[rightpart] (0.08,0) -- (4.2,0);
  \draw[marker] (0,-0.78) -- (0,0.78);

  \foreach \x in {-3.35,-2.55,-1.75,-0.95,-0.28} {
    \node[qpoint, blue!65!black] at (\x,0) {};
  }
  \foreach \x in {0.35,1.1,1.85,2.72,3.55} {
    \node[qpoint, orange!80!black] at (\x,0) {};
  }

  \node[blue!55!black, above] at (-2.15,0.25) {$A=\alpha$};
  \node[orange!65!black, above] at (2.15,0.25) {$B=\Q\setminus\alpha$};
  \node[gray!55!black, below] at (0,-0.16) {cut};
\end{tikzpicture}

  \caption{A Dedekind cut splits $\Q$ into $\alpha$ and $\Q\setminus\alpha$, with no last rational in $\alpha$.}
  \label{fig:dedekind-cut-left-right-parts}
\end{figure}
 
\begin{remark} \label{rem:left-cut-versus-two-sided-cut} \hlt{Left Cuts Versus Two-Sided Cuts}\\
Two-sided notation writes a cut as $A|B$, with $A=\alpha$ and $B=\Q\setminus\alpha$; Definition \ref{def:dedekind-cuts-and-order} records only $A$.\\
Thus ``no largest element'' is the same left-side condition.
\end{remark}

\begin{definition} \label{def:cut-arithmetic} \hlt{Arithmetic of Cuts}\\
For cuts $\alpha,\beta$, define $\alpha+\beta\defeq\{r+s:r\in\alpha,\ s\in\beta\}$ and $0^*\defeq\{q\in\Q:q<0\}$.\\
Let candidate additive inverse of $\alpha$ be $-\alpha\defeq\{p\in\Q:\exists r\in\Q,\ r>0,\ -p-r\notin\alpha\}$. Let $\mathscr{D}^+\defeq\{\alpha\in\mathscr{D}:\alpha>0^*\}$.\\
For $\alpha,\beta\in\mathscr{D}^+$, define $\alpha\beta\defeq\{p\in\Q:\exists r\in\alpha,\ \exists s\in\beta,\ r>0,\ s>0,\ p\leq rs\}$ and $1^*\defeq\{q\in\Q:q<1\}$.\\
Extend multiplication to all cuts by setting $\alpha0^*=0^*\alpha=0^*$ and, for nonzero signs,
\begin{align}
\alpha\beta \defeq
\begin{cases}
(-\alpha)(-\beta), & \alpha<0^*,\ \beta<0^*,\\
-\bigl((-\alpha)\beta\bigr), & \alpha<0^*,\ \beta>0^*,\\
-\bigl(\alpha(-\beta)\bigr), & \alpha>0^*,\ \beta<0^*.
\end{cases} \nonumber
\end{align}
\hlt{Motivation}: arithmetic uses rational approximations inside cuts; signs come after positive products.
\end{definition}

\begin{definition} \label{def:rational-cuts} \hlt{Rational Cuts}\\
For $r\in\Q$, define the \hlt{rational cut}, $r^*\defeq\{p\in\Q:p<r\}$.\\
The map $r\mapsto r^*$ is intended to identify $\Q$ with a subfield of the cut field.
\end{definition}

\begin{theorem} \label{thm:real-field-existence-lub} \hlt{Existence of the Real Field}\\
There exists an ordered field $\R$ with the least-upper-bound property. Moreover, $\R$ contains $\Q$ as a subfield.
\begin{proof}
Construct $\R$ from $\Q$ using Definitions \ref{def:dedekind-cuts-and-order}--\ref{def:rational-cuts}.
\begin{enumerate}[label=(\roman*)]
\setlength{\itemsep}{0pt}
\item \hlt{Order.}
The relation $\alpha<\beta$ is transitive, and at most one of $\alpha<\beta$, $\alpha=\beta$, and $\beta<\alpha$ can hold. For trichotomy, suppose $\alpha\not\subseteq\beta$ and $\alpha\neq\beta$. Choose $p\in\alpha$ with $p\notin\beta$. If $q\in\beta$, then $q<p$, since $p\notin\beta$ and $\beta$ is downward closed. Hence $q\in\alpha$, so $\beta\subseteq\alpha$. Since $\alpha\neq\beta$, $\beta<\alpha$. Thus $\mathscr{D}$ is ordered.\\
\hlt{Commentary}: One missing rational decides containment.

\item \hlt{Least Upper Bounds.}
Let $A\subseteq\mathscr{D}$ be nonempty and bounded above by a cut $\beta$. Define $\gamma\defeq\bigcup\limits_{\alpha\in A}\alpha$.\\
Since some $\alpha_0\in A$ is nonempty, $\gamma$ is nonempty. Since every $\alpha\in A$ lies in $\beta$, $\gamma\subseteq\beta$, so $\gamma\neq\Q$. If $p\in\gamma$ and $q<p$, then $p\in\alpha_1$ for some $\alpha_1\in A$, hence $q\in\alpha_1\subseteq\gamma$. If $p\in\gamma$, choose $\alpha_1\in A$ with $p\in\alpha_1$ and then choose $r\in\alpha_1$ with $p<r$; then $r\in\gamma$. Thus $\gamma$ is a cut.\\
Every $\alpha\in A$ satisfies $\alpha\leq\gamma$. If $\delta<\gamma$, choose $s\in\gamma\setminus\delta$. Then $s\in\alpha$ for some $\alpha\in A$, so $\delta<\alpha$ and $\delta$ is not an upper bound of $A$. Hence $\gamma=\sup A$.\\
\hlt{Motivation}: Supremum is union.\\
\hlt{Commentary}: Bounded unions stay proper and give the smallest upper bound.

\item \hlt{Addition.}
The set $\alpha+\beta$ is a cut. It is nonempty by choosing one element from each summand.\\
If $r'\notin\alpha$ and $s'\notin\beta$, then every $r\in\alpha$ and $s\in\beta$ satisfies $r<r'$ and $s<s'$, so $r+s<r'+s'$; hence $r'+s'\notin\alpha+\beta$, and $\alpha+\beta$ is proper. If $p=r+s\in\alpha+\beta$ and $q<p$, then $q=r+(s+q-p)$ with $s+q-p<s$, so $s+q-p\in\beta$ and $q\in\alpha+\beta$. If $p=r+s$, choose $r_1\in\alpha$ with $r<r_1$; then $p<r_1+s\in\alpha+\beta$.\\
Thus $\alpha+\beta$ is a cut. Commutativity and associativity follow by rewriting representatives in $\Q$.\\
The cut $0^*$ is the additive identity: $\alpha+0^*\subseteq\alpha$ by downward closure, while if $p\in\alpha$, choose $r\in\alpha$ with $p<r$ and write $p=r+(p-r)$ with $p-r<0$.\\
The inverse $-\alpha$ is a cut, and $\alpha+(-\alpha)=0^*$. The inclusion $\alpha+(-\alpha)\subseteq0^*$ follows because $s\in-\alpha$ forces $-s\notin\alpha$, so any $r\in\alpha$ satisfies $r<-s$. Conversely, if $v<0$, put $w=-v/2>0$ and choose an integer $n$ with $nw\in\alpha$ but $(n+1)w\notin\alpha$; then $p=-(n+2)w$ lies in $-\alpha$ and $v=nw+p\in\alpha+(-\alpha)$.\\
\hlt{Commentary}: Addition inherits its laws from $\Q$ once closure and inverses are checked.

\item \hlt{Order Compatibility with Addition.}
The additive laws make $\mathscr{D}$ an additive group.\\
If $\beta<\gamma$, then $\alpha+\beta\subseteq\alpha+\gamma$. Equality would imply $\beta=\gamma$ after adding $-\alpha$ to both sides, so $\alpha+\beta<\alpha+\gamma$. Consequently, $\alpha>0^*$ if and only if $-\alpha<0^*$.\\
\hlt{Commentary}: Cancellation turns inclusion into strict inclusion.

\item \hlt{Positive Multiplication.}
If $\alpha,\beta\in\mathscr{D}^+$, then $\alpha\beta$ is a cut. Since each positive cut contains some positive rational, $\alpha\beta$ is nonempty. Choose positive $r'\notin\alpha$ and $s'\notin\beta$; then every positive $r\in\alpha$ and $s\in\beta$ satisfies $rs<r's'$, so $r's'\notin\alpha\beta$, and the product is proper. Downward closure is built into the condition $p\leq rs$. If $p\leq rs$ with $r,s>0$, then either choose $q$ with $p<q<rs$, or, when needed, choose $r_1\in\alpha$ with $r<r_1$ and take $q$ with $p<q\leq r_1s$; hence the product has no largest element.\\
Commutativity follows from $rs=sr$. Associativity and distributivity are proved by unfolding membership and using the rational identities $(rs)t=r(st)$ and $r(s+t)=rs+rt$, with a positive rational slack inserted whenever a representative must be made strictly larger inside a cut. The identity is $1^*$: if $p\leq rs$ with $r\in\alpha$, $s<1$, and $r,s>0$, then $p<r$, so $p\in\alpha$; conversely, if $p\in\alpha$, choose $r\in\alpha$ with $p<r$ and then choose $s\in\Q$ with $0<s<1$ and $p\leq rs$. In particular, if $\alpha>0^*$ and $\beta>0^*$, then $\alpha\beta>0^*$.\\
\hlt{Commentary}: Positivity avoids sign reversals.

\item \hlt{Signed Multiplication.}
The definitions extend multiplication to all cuts by reducing to positive factors and additive inverses. Since $\delta=-(-\delta)$ and the sign rules for addition have already been proved, commutativity, associativity, and inverses reduce to the positive case.\\
For distributivity, split into signs of $\alpha,\beta,\gamma,\beta+\gamma$. For example, if $\alpha>0^*$, $\beta<0^*$, and $\beta+\gamma>0^*$, then $\gamma=(\beta+\gamma)+(-\beta)$, so positive distributivity gives $\alpha\gamma=\alpha(\beta+\gamma)+\alpha(-\beta)$. Because $\alpha(-\beta)=-(\alpha\beta)$, adding $\alpha\beta$ to both sides gives $\alpha\beta+\alpha\gamma=\alpha(\beta+\gamma)$. The other sign cases are identical reductions after replacing negative terms by their additive inverses.\\
\hlt{Commentary}: Distributivity reduces to positive cases plus signs.

\item \hlt{Embedding $\Q$.}
Each $r^*$ is a cut. For $r,s\in\Q$, $r^*+s^*=(r+s)^*$, $r^*s^*=(rs)^*$, and $r^*<s^*\Longleftrightarrow r<s$.\\
For addition, if $p=u+v$ with $u<r$ and $v<s$, then $p<r+s$. Conversely, if $p<r+s$, choose $t>0$ with $2t=r+s-p$ and put $r'=r-t$, $s'=s-t$; then $r'<r$, $s'<s$, and $p=r'+s'$. For positive products, if $p\leq uv$ with $0<u<r$ and $0<v<s$, then $p<rs$. Conversely, if $0<p<rs$, choose positive $u<r$ and $v<s$ close enough to $r$ and $s$ that $p\leq uv$; if $p\leq0$, use any positive $u<r$, $v<s$. The signed product cases follow from the definitions by additive inverses. The order equivalence follows because $r<s$ puts $r$ in $s^*$ but not in $r^*$, while $r^*<s^*$ gives some $p<s$ with $p\notin r^*$, hence $r\leq p<s$.\\
\hlt{Commentary}: The map $r\mapsto r^*$ preserves arithmetic and order.
\end{enumerate}
Identify $\Q$ with the subfield $\Q^*\defeq\{r^*:r\in\Q\}\subseteq\mathscr{D}$, and write $\R$ for $\mathscr{D}$. Items (i), (iii)--(vi) give an ordered field; item (ii) gives the least-upper-bound property; item (vii) embeds $\Q$ as a subfield.\\
Any two ordered fields with the least-upper-bound property are isomorphic, so this construction characterizes the real field up to ordered-field structure.
\end{proof}
\end{theorem}

\begin{perspective} \label{per:number-systems-from-representatives} \hlt{Number Systems from Representatives}\\
Choose representatives $X$, an equivalence relation $\sim$, and classes $[x]\in X/{\sim}$: $\Z:(a,b)\mapsto[(a,b)]$, $\Q:(a,b)$ with $b\neq0\mapsto[(a,b)]$, and $\R$ via cuts or Cauchy data. Prove each operation or order is well defined, e.g. $x\sim x'$ and $y\sim y'$ imply $x\star y\sim x'\star y'$ and $x<y\Leftrightarrow x'<y'$; then use $X/{\sim}$ as the number system.\\
\hlt{Key Insight}: Laws must descend to classes.
\end{perspective}

\begin{theorem} \label{thm:archimedean-and-rational-density} \hlt{Archimedean Property and Density of $\Q$}
\begin{enumerate}[label=(\roman*)]
\setlength{\itemsep}{0pt}
\item If $x,y\in\R$ and $x>0$, then there is a positive integer $n$ such that $nx>y$.
\item If $x,y\in\R$ and $x<y$, then there exists $p\in\Q$ such that $x<p<y$.
\end{enumerate}
\begin{proof}
\begin{enumerate}[label=(\roman*)]
\setlength{\itemsep}{0pt}
\item Let $A=\{nx:n\in\N\}$. If no $nx$ exceeded $y$, then $A$ would be bounded above.\\
By Theorem \ref{thm:real-field-existence-lub}, let $\alpha=\sup A$. Since $x>0$, $\alpha-x<\alpha$, so $\alpha-x$ is not an upper bound.\\
Hence $\alpha-x<mx$ for some $m\in\N$, and $\alpha<(m+1)x\in A$, a contradiction.
\item Since $y-x>0$, part (i) gives $n\in\N$ with $n(y-x)>1$. Apply part (i) again to get $m_1,m_2\in\N$ with $m_1>nx$ and $m_2>-nx$. Then some $m\in\Z$ satisfies $m-1\leq nx<m$.\\
Thus $nx<m\leq nx+1<ny$, so $x<m/n<y$ and $m/n\in\Q$.
\end{enumerate}
\hlt{Structural Commentary}: Completeness makes $\{nx:n\in\N\}$ unbounded.\\
\hlt{Key Insight}: Archimedean unboundedness gives rational density.
\end{proof}
\end{theorem}

\begin{figure}[H]
  \centering
  \begin{tikzpicture}[
  >=Stealth,
  every node/.style={font=\small},
  axis/.style={gray!70, line width=0.55pt},
  tick/.style={gray!70, line width=0.55pt},
  point/.style={circle, fill=blue!70!black, inner sep=1.55pt},
  stepbar/.style={orange!75!black, line width=0.85pt}
]
  \draw[axis] (-0.4,0) -- (8.7,0);
  \foreach \x in {0,1,2,3,4,5,6,7,8} {
    \draw[tick] (\x,-0.1) -- (\x,0.1);
  }

  \node[below] at (0,-0.1) {$0$};
  \node[above] at (1,0.1) {$\frac{1}{n}$};
  \node[above] at (2,0.1) {$\frac{2}{n}$};
  \node[above] at (3,0.1) {$\frac{3}{n}$};
  \node[above] at (4,0.1) {$\cdots$};
  \node[above] at (5,0.18) {$\frac{m-1}{n}$};
  \node[above] at (6,0.18) {$\frac{m}{n}$};

  \node[point] at (5.42,0) {};
  \node[point] at (6.9,0) {};
  \node[below] at (5.42,-0.1) {$a$};
  \node[below] at (6.9,-0.1) {$b$};

  \draw[stepbar] (5,-0.48) -- (6,-0.48);
  \draw[stepbar] (5,-0.58) -- (5,-0.38);
  \draw[stepbar] (6,-0.58) -- (6,-0.38);
  \node[below, orange!75!black] at (5.5,-0.58) {$1/n<b-a$};
\end{tikzpicture}

  \caption{Choose $n$ so the spacing $1/n$ is smaller than $b-a$; then one grid point $m/n$ falls inside $(a,b)$.}
  \label{fig:rational-density-step-size}
\end{figure}

\begin{theorem} \label{thm:nth-roots-positive-reals} \hlt{Existence and Uniqueness of Positive $n$th Roots}\\
For every $x\in\R$ with $x>0$ and every positive integer $n$, $\exists$ a unique positive real number $y$ such that $y^n=x$.
\begin{proof}
\begin{enumerate}[label=(\roman*)]
\setlength{\itemsep}{0pt}
\item \hlt{Uniqueness.} If $0<y_1<y_2$, then $y_1^n<y_2^n$.
\item \hlt{Existence.} Let $E\defeq\{t\in\R:t>0,\ t^n<x\}$. $E\neq\varnothing$: for $t=x/(1+x)$, $0<t<1$ and $t^n\leq t<x$.\\
Also $E$ is bounded above by $1+x$, since $t>1+x$ implies $t^n\geq t>x$. By Theorem \ref{thm:real-field-existence-lub}, let $y=\sup E$.\\
For $0<a<b$, $b^n-a^n=(b-a)(b^{n-1}+b^{n-2}a+\cdots+a^{n-1})<(b-a)nb^{n-1}$.\\
If $y^n<x$, choose $h$ with $0<h<1$ and $h<\frac{x-y^n}{n(y+1)^{n-1}}$.\\
Then $(y+h)^n-y^n<hn(y+h)^{n-1}<hn(y+1)^{n-1}<x-y^n$.\\
Thus $y+h\in E$, contradicting that $y$ is an upper bound. If $y^n>x$, set $k\defeq\frac{y^n-x}{ny^{n-1}}$. Then $0<k<y$.\\
For $t\geq y-k$, $y^n-t^n\leq y^n-(y-k)^n<kny^{n-1}=y^n-x$, so $t^n>x$ and $t\notin E$.\\
Hence $y-k$ is an upper bound for $E$, contradicting $y=\sup E$. Therefore $y^n=x$.
\end{enumerate}
\hlt{Structural Commentary}: The root is $\sup\{t>0:t^n<x\}$.\\
\hlt{Key Insight}: Completeness supplies the missing root.
\end{proof}
\end{theorem}

\begin{corollary} \label{cor:rationals-and-irrationals-dense} \hlt{Density of Rational and Irrational Points}\\
If $x<y$ in $\R$, then there exist $r\in\Q$ and $s\in\R\setminus\Q$ such that $x<r<y$ and $x<s<y$.
\begin{proof}
The rational point $r$ is supplied by Theorem \ref{thm:archimedean-and-rational-density}. Let $\eta=\sqrt{2}$.\\
Choose $n\in\N$ with $0<\eta/n<y-x$, and choose $q\in\Q$ with $x-\eta/n<q<y-\eta/n$.\\
Then $s\defeq q+\eta/n$ lies in $(x,y)$ and is irrational, since $s\in\Q$ would imply $\eta=n(s-q)\in\Q$.\\
\hlt{Structural Commentary}: Irrational density is rational density shifted by a small irrational amount.\\
\hlt{Key Insight}: Once one irrational is known, every interval contains irrational translates of rational points.
\end{proof}
\end{corollary}

\begin{corollary} \label{cor:nth-root-product-positive-reals} \hlt{Product Rule for Positive Roots}\\
If $a,b>0$ and $n$ is a positive integer, then $(ab)^{1/n}=a^{1/n}b^{1/n}$.
\begin{proof}
Let $\alpha=a^{1/n}$ and $\beta=b^{1/n}$. Then $ab=\alpha^n\beta^n=(\alpha\beta)^n$ by commutativity of multiplication.\\
The uniqueness part of Theorem \ref{thm:nth-roots-positive-reals} gives the identity.\\
\hlt{Structural Commentary}: Use uniqueness after checking $n$th powers.\\
\hlt{Key Insight}: Root identities are uniqueness arguments.
\end{proof}
\end{corollary}

\begin{definition} \label{def:extended-real-number-system} \hlt{Extended Real Number System}\\
The \hlt{extended real number system} is $\overline{\R}\defeq\R\cup\{-\infty,+\infty\}$.\\
The order on $\R$ is preserved, and $-\infty<x<+\infty$, $(x\in\R)$.\\
Every subset of $\R$ has a supremum and an infimum in $\overline{\R}$ after the conventions $\sup\emptyset=-\infty$, $\inf\emptyset=+\infty$, $\sup E=+\infty$ when $E$ is unbounded above, and $\inf E=-\infty$ when $E$ is unbounded below.
\end{definition}

\begin{remark} \label{rem:extended-real-arithmetic-conventions} \hlt{Extended Arithmetic Conventions}\\
$\overline{\R}$ is not a field. For $x\in\R$:
\begin{enumerate}[label=(\roman*)]
\setlength{\itemsep}{0pt}
\item $\infty+\infty=\infty$ and $-\infty-\infty=-\infty$.
\item $x+\infty=\infty$ and $x-\infty=-\infty$.
\item If $x>0$, then $x\cdot\infty=\infty$ and $x\cdot(-\infty)=-\infty$; if $x<0$, the signs reverse.
\end{enumerate}
Expressions $\infty-\infty$ and $0\cdot\infty$ are undefined.
\end{remark}

\begin{remark} \label{rem:unbounded-real-intervals} \hlt{Unbounded Intervals}\\
For $a,b\in\R$:
\begin{enumerate}[label=(\roman*)]
\setlength{\itemsep}{0pt}
\item $(a,\infty)\defeq\{x\in\R:a<x\}$ and $[a,\infty)\defeq\{x\in\R:a\leq x\}$.
\item $(-\infty,b)\defeq\{x\in\R:x<b\}$ and $(-\infty,b]\defeq\{x\in\R:x\leq b\}$.
\end{enumerate}
Also $(-\infty,\infty)\defeq\R$.
\end{remark}
